\documentclass[a4paper, landscape]{article}
\usepackage[8pt]{extsizes}
\usepackage[utf8]{inputenc}
\usepackage[russian]{babel}
\usepackage{amsmath, amssymb, amsthm, mathtools}
\usepackage[margin=0.4cm]{geometry}
\usepackage{multicol}
\usepackage{enumitem}
\usepackage{sectsty}

\sectionfont{\normalsize\bfseries\scshape\vspace{-2ex}\hrule\vspace{1ex}}
\subsectionfont{\small\bfseries\vspace{-2ex}}
\setlength{\parindent}{0pt}
\setlength{\parskip}{1pt}
\setlength{\abovedisplayskip}{1pt}
\setlength{\belowdisplayskip}{1pt}
\setlist{nosep, leftmargin=1.2em}

\newcommand{\E}{\mathbb{E}}
\renewcommand{\P}{\mathbb{P}}
\newcommand{\N}{\mathcal{N}}
\DeclareMathOperator{\Var}{Var}
\DeclareMathOperator{\Cov}{Cov}
\newcommand{\PhiFunc}{\Phi} % Функция Лапласа

\begin{document}
\begin{multicols*}{4}

\section*{1. Асимптотика и ЛБК}
\textbf{Леммы Бореля-Кантелли (ЛБК):}
\begin{enumerate}
    \item $\sum \P(A_n) < \infty \implies \P(A_n \text{ б.ч.}) = 0$.
    \item $\sum \P(A_n) = \infty$, $A_n$ незав. $\implies \P(A_n \text{ б.ч.}) = 1$.
\end{enumerate}
\textbf{Контрпример ($\xrightarrow{\P} \not\implies \xrightarrow{п.н.}$):} «Бегущий индикатор» на $[0,1]$ (отрезки длины $1/n$). Индикатор $\xrightarrow{\P} 0$, но любая точка покрывается бесконечно часто (т.к. $\sum 1/n \!=\! \infty$, ЛБК-2), п.н. предела нет. Сходимость есть, только если ряд вероятностей сходится!

\textbf{Алгоритм поиска $C = \varlimsup \frac{\xi_n - a_n}{b_n}$:}
\begin{enumerate}
    \item \textbf{Порог:} Вводим $\varepsilon>0$, $A_n^\pm = \{ \xi_n > a_n + b_n(C \pm \varepsilon) \}$. Заменяем $\P(A_n)$ на асимптотику хвоста.
    \item \textbf{Ряды (Шкала Бертрана):} Обобщ. гарм. ряд $\sum \frac{(\ln n)^\beta}{n(\ln n)^C} \propto \sum \frac{1}{n(\ln n)^q}$. Сходится при $q>1$, расходится при $q \le 1$.
    \item \textbf{Синтез:} Подбираем $C$: для $A_n^+$ ряд сходится (ЛБК-1 $\implies \varlimsup \le C$), для $A_n^-$ расходится (ЛБК-2 $\implies \varlimsup \ge C$).
\end{enumerate}

\textbf{Асимптотика хвостов $\P(\xi > x)$ при $x \to \infty$:}
\begin{itemize}
    \item $\N(0,1): \sim \frac{1}{\sqrt{2\pi} x} e^{-x^2/2}$ (убывает как $x^{-1}$)
    \item $\Gamma(\alpha, \lambda): \sim \frac{\lambda^{\alpha-1}}{\Gamma(\alpha)} x^{\alpha-1} e^{-\lambda x}$
    \item $\text{Exp}(\lambda): = e^{-\lambda x}$ (точное равенство)
    \item Тяж. хвосты (Коши): $\sim C x^{-\alpha}$. Нормировка $b_n=n^{1/\alpha}$.
\end{itemize}

\textbf{Мастер-теоремы для экстремумов (Анализ через ЛБК и ряды Бертрана)} \\

\begin{enumerate}[leftmargin=*, nosep, itemsep=4pt, label=\textbf{\arabic*.}]
    \item \textbf{Абсолютная формула (Экспоненциального типа):} \\
    Пусть асимптотика хвоста имеет вид $\P(\xi_n > x) \sim K \cdot x^\beta e^{-\lambda x}$. \\
    При \textbf{стандартной нормировке} $a_n = \frac{1}{\lambda}\ln n$ и $b_n = \frac{1}{\lambda}\ln \ln n$, критический порог вычисляется по формуле: 
    $\mathbf{C = \beta + 1}$. 
    
    \textit{Внимание к знаменателю!} Если в условии нормировка задана как $b_n = \ln \ln n$ (без $\frac{1}{\lambda}$), уравнение баланса степеней имеет вид $\lambda C - \beta = 1 \implies \mathbf{C = \frac{\beta + 1}{\lambda}}$.

    \item \textbf{Экспоненциальное $\text{Exp}(\lambda)$:} \\
    Асимптотика точная: $\P(\xi_n > x) = e^{-\lambda x}$. Степень перед экспонентой $\beta = 0$. \\
    $\mathbf{C = 1}$.

    \item \textbf{Гамма-распределение $\Gamma(\alpha, \lambda)$ (вкл. Эрланга и $\chi^2$):} \\
    Асимптотика: $\P(\xi_n > x) \sim \frac{\lambda^{\alpha-1}}{\Gamma(\alpha)} x^{\alpha-1} e^{-\lambda x}$. Степень $\beta = \alpha - 1$. \\
     $\mathbf{C = \alpha}$.

    \item \textbf{Нормальное $\mathcal{N}(0, 1)$:} \\
    Асимптотика (отношение Миллса): $\P(\xi_n > x) \sim \frac{1}{\sqrt{2\pi}} x^{-1} e^{-x^2/2}$. \\
    Нормировка: $a_n = \sqrt{2\ln n}$. Степень перед экспонентой $\beta = -1$. \\
     $\mathbf{C = -1/2}$.

    \item \textbf{Тяжелые хвосты (Парето, Коши, Фреше):} \\
    Асимптотика убывает степенным образом: $\P(\xi_n > x) \sim K \cdot x^{-\alpha}$. \\
    Экспоненциальная мастер-теорема \textbf{не работает}! Для таких распределений нормирующая последовательность всегда имеет степенной вид $\mathbf{b_n = n^{1/\alpha}}$, а не логарифмический.
\end{enumerate}

\section*{2. Сходимости и Дельта-метод}
\textbf{Топология сходимостей:}
\begin{itemize}
    \item $\xi_n \xrightarrow{d} C \iff \xi_n \xrightarrow{\P} C$ (\textbf{только} если предел --- константа!).
    \item ЦПТ: $\eta_n = \sqrt{n}(\overline{X} - \mu)$ сходится по $d$, но \textbf{не имеет} предела по $\P$ ($\eta_n - \eta_m$ не стремится к $0$).
    \item Вектора: $\xi_n \xrightarrow{d} \xi, \eta_n \xrightarrow{d} \eta \not\Rightarrow (\xi_n, \eta_n) \xrightarrow{d} (\xi, \eta)$. Нужна независимость компонент или многомерная ЦПТ.
\end{itemize}

\textbf{Алгоритм Дельта-метода (Сходимость по $d$):}
\begin{enumerate}
    \item \textbf{База:} Найти параметры $X_1$: $\mu = \E X_1$, $\sigma^2 = \Var X_1$.
    \item \textbf{ЦПТ:} Записать базу: $\sqrt{n}(\overline{X} - \mu) \xrightarrow{d} \N(0, \sigma^2)$.
    \item \textbf{Функция:} Выделить $g(t)$. Убедиться, что вычитаемое $C = g(\mu)$.
    \item \textbf{Дисперсия:} Найти $g'(\mu)$. Асимпт. дисперсия $\sigma_{as}^2 = [g'(\mu)]^2 \sigma^2$. Предел: $\N(0, \sigma_{as}^2)$.
\end{enumerate}

\textbf{ВЫРОЖДЕНИЕ (2-й порядок, Теорема о $\chi^2$):}
Если $g'(\mu)=0$, но $g''(\mu)\ne 0$, классика дает $\sigma_{as}^2=0$ ($\xrightarrow{\P} 0$). Домножаем ряд Тейлора на $n$ (а не на $\sqrt{n}$):
$ n(g(\overline{X}) - g(\mu)) \xrightarrow{d} \frac{\sigma^2 g''(\mu)}{2} \chi^2_1 $

\textit{Предел — гамма-распределение (пропорц. $\chi^2$), а не норм!}

\textbf{Многомерный Дельта-метод:} $\vec{W}_i=(Y_i, Z_i)^T$ --- выборка, $\E[\vec{W}_1]=\vec{\mu}$, $\Sigma = \Cov(\vec{W}_1)$. $\sigma_{as}^2 = \nabla g(\vec{\mu})^T \Sigma \nabla g(\vec{\mu})$.
\textit{Синтез для $g(y,z) = z/y^k$:} $\nabla g = \left(-\frac{k\mu_z}{\mu_y^{k+1}}, \frac{1}{\mu_y^k}\right)^T$.
$$ \sigma^2_{as} \!=\! \frac{k^2 \mu_z^2}{\mu_y^{2k+2}} \Var(Y) - \frac{2k \mu_z}{\mu_y^{2k+1}} \Cov(Y, Z) + \frac{1}{\mu_y^{2k}} \Var(Z) $$

\section*{3. Гауссовские векторы}
\textbf{ВНИМАНИЕ (Сдвиг средних):} Проекции и формула Вика применимы \textbf{ТОЛЬКО} к центрированным векторам ($\vec{\mu}=\vec{0}$). Иначе сначала замена: $X_0 = X-\mu_X$, $Y_0=Y-\mu_Y$.

\textbf{Формула Шеппарда (Вероятности ортантов):} 
Для центрированного $(X,Y)$ с корреляцией $\rho$:
$\P(X \!>\! 0, Y \!>\! 0) \!=\! \frac{1}{4} + \frac{1}{2\pi}\arcsin(\rho)$, \ $\P(XY \!<\! 0) \!=\! \frac{1}{\pi}\arccos(\rho)$.

\textbf{Многомерная ХФ (Экспонента от вектора):}
Для $\vec{\xi} \sim \N(\vec{\mu}, \Sigma)$ интеграл берется аналитически:
$$ \E[\exp(\vec{t}^T \vec{\xi})] = \exp\left( \vec{t}^T \vec{\mu} + \frac{1}{2} \vec{t}^T \Sigma \vec{t} \right) $$
\textit{Пример:} $\E[e^{X+Y+Z}] \implies \vec{t}=(1,1,1)^T$, ответ --- экспонента от полусуммы всех элементов $\Sigma$.

\textbf{Формула Вика (Теорема Иссерлиса):} Для $\N(\vec{0}, \Sigma)$: $\E[\text{нечетн.}]=0$. $\E[X_1 X_2 X_3 X_4] = \E[X_1 X_2]\E[X_3 X_4] + \E[X_1 X_3]\E[X_2 X_4] + \E[X_1 X_4]\E[X_2 X_3]$.

\textbf{Алгоритм вычисления нелинейного $\E[\cos(YZ)|X]$:}
\begin{enumerate}
    \item \textbf{Проекции:} Представить $Y \!=\! aX+\tilde{Y}$, $Z \!=\! bX+\tilde{Z}$.
    $a \!=\! \frac{\Cov(Y,X)}{\Var(X)}$, $\Var(\tilde{Y}) \!=\! \Var(Y) - \frac{\Cov^2(X,Y)}{\Var(X)}$.
    Вычислить $\Cov(\tilde{Y},\tilde{Z}) \!=\! \Cov(Y,Z) - ab\Var(X)$. Если 0 $\implies \tilde{Y} \perp\!\!\perp \tilde{Z}$. Зафиксируем $X=x \implies Y_x, Z_x$ независимы.
    \item \textbf{Комплексификация:} Интеграл косинуса не берется. Формула Эйлера: $\cos \alpha = \text{Re}(e^{i\alpha})$. Ищем $\text{Re}(\E[e^{iY_x Z_x}])$.
    \item \textbf{Итерированное $\E$:} По свойству башни обусловим по $Z_x$. Внутреннее $\E$ --- это ХФ нормальной $Y_x$:
    $\E[e^{i Z_x Y_x}|Z_x] = \exp(i\mu_{Y_x} Z_x - \frac{1}{2}\sigma_{Y_x}^2 Z_x^2)$.
    \item \textbf{Интеграл Гаусса:} Внешнее $\E$ по $Z_x \sim \N(\mu_{Z_x}, \sigma_{Z_x}^2)$ (с константами $c=i\mu_{Y_x}, d=-\frac{1}{2}\sigma_{Y_x}^2$):

    $ \E[e^{c\eta + d\eta^2}] = \frac{1}{\sqrt{1-2d\sigma^2}}\exp\left( \frac{c^2\sigma^2+2c\mu+2d\mu^2}{2(1-2d\sigma^2)} \right) $
\end{enumerate}

\section*{4. Смешанные распр. и ожидания}
\textbf{Тонелли и Свойство Башни ($X \perp\!\!\perp Y$):} Задача $\E[g(X,Y)]$, где $X$ со скачками, $Y$ непрерывна. $\E[g(X,Y)] = \E[\E[g(X,Y)|X]] = \E[\psi(X)]$. Сводим двумерный интеграл к двум одномерным:
\begin{enumerate}
    \item \textbf{Внутреннее $\E$:} По непрерывной $Y$. Фиксируем $X=x$. $\psi(x) = \E[g(x,Y)] = \int g(x,y)p_Y(y)dy$. \textit{Трюк:} Для $Y^X$ интегрируем $Y$ первой, избегая трансцендентных $e^{X\ln y}$!
    \textbf{БЫСТРАЯ ФОРМУЛА:} Если $Y \!\sim\! U(a,b)$, то $\psi(x) = \E[Y^x] = \int_a^b y^x \frac{1}{b-a} dy = \frac{b^{x+1}-a^{x+1}}{(b-a)(x+1)}$.
    \item \textbf{Дельта-функция для $X$:} Обобщенная плотность $p_X^*(x) = F'_{ac}(x) + \sum p_k \delta(x-x_k)$. Атомы: точки разрыва $x_k$, их веса $p_k = F_X(x_k) - F_X(x_k-0)$.
    \item \textbf{Синтез:} Итоговое $\E$ по обобщенной плотности:
    $$ \E[g(X,Y)] = \int_{\mathbb{R}\setminus D} \psi(x) F'_{ac}(x) dx + \sum_{k} \psi(x_k) p_k $$
\end{enumerate}
\textbf{ОПАСНОСТЬ (Ложная нормировка):} Непрерывная часть вне разрывов $p_{ac}(x) = F'_X(x)$. Строго: $\int p_{ac}dx = 1 - \sum p_k < 1$. Применять табличные формулы $\E$ только к $p_{ac}$ в отрыве от $p_k$ недопустимо!

\section*{5. Условные плотности и Якобиан}
\textbf{Алгоритм Якобиана (поиск $p_{X|Z}$, где $Z=f(X,Y)$):}
\begin{enumerate}
    \item \textbf{Совместная:} Если $X \perp\!\!\perp Y \implies p_{X,Y} = p_X(x)p_Y(y)$.
    \item \textbf{Замена:} Вводим $U\!=\!X, Y\!=\!h(X,Z)$. Вычисляем якобиан $J = \partial h(x,z) / \partial z$. Плотность: $p_{X,Z} = p_{X,Y}(x, h(x,z)) \cdot |J|$.
    \item \textbf{ДИНАМИЧЕСКИЕ ГРАНИЦЫ (Критично!):} Подставляем $y=h(x,z)$ в индикаторы исходной плотности $I_{\{y \in D_y\}}$. Разрешаем систему относительно $x \implies A(z) < x < B(z)$.
    \textit{Пример:} $x \!\in\! (0,1), xz \!\in\! (0,1) \implies 0 < x < \min(1, 1/z)$.
    \item \textbf{Маргинализация:} Интегрируем $p_{X,Z}$ по $x$. Из-за $\min$ интеграл $p_Z(z)$ распадается на ветви (напр. $z \le 1$ и $z > 1$).
    \item \textbf{Условная:} $p_{X|Z} \!=\! p_{X,Z}/p_Z$. $\E[g|Z=z] \!=\! \int g(\dots) p_{X|Z}dx$.
\end{enumerate}

\textbf{Трюк измеримости (Вынесение множителя):}
Относительно условия $Z$, сама $Z$ — константа.
$\E[XY | Z=Y/X] = \E[X(XZ)|Z] = Z \cdot \E[X^2|Z]$. (Ищем только $\E[X^2|Z]$!)

\textbf{Комбинаторная симметрия:} Если $X_1 \dots X_n$ н.о.р., а функция условия $S = \sum X_i$ симметрична, то все $\E[X_k|S]$ равны.
$\sum \E[X_k|S] = \E[S|S] = S \implies \E[X_k|S] = S/n$. Без интегралов!

\textbf{Немонотонные условия:} Если биекции нет ($Z=X^2$).
$X \sim U(-a,a)$, плотность четная. Усл. распр. дискретно: $\P(X\!=\!\sqrt{z}|z) = \P(X\!=\!-\sqrt{z}|z) = 1/2 \implies \E[X|X^2=z] = 0$.

\section*{Характеристические функции}
$\varphi(t) = \E[e^{it\xi}]$. Свойства: $\varphi(0)\!=\!1$, $|\varphi(t)| \!\le\! 1$, $\varphi_{aX+b}(t) \!=\! e^{itb}\varphi_X(at)$. Моменты: $\E[X^k] = \varphi^{(k)}(0) / i^k$.
$\xi \perp\!\!\perp \eta \implies \varphi_{\xi+\eta} = \varphi_\xi \varphi_\eta$. (Обратное неверно).
\textbf{Решетчатость:} $|\varphi(t_0)|=1$ для $t_0 \ne 0 \iff$ распр. дискретное.
\textbf{Таблица ХФ:}
\begin{itemize}
    \item $\N(\mu,\sigma^2): \exp(i\mu t - \frac{1}{2}\sigma^2 t^2)$
    \item $\text{Po}(\lambda): \exp(\lambda(e^{it}-1))$
    \item $\text{Bin}(n,p): (1-p+pe^{it})^n$
    \item $\text{Exp}(\lambda): \lambda / (\lambda-it)$
    \item $U(a,b): (e^{itb}-e^{ita}) / (it(b-a))$
    \item Лаплас $L(0,b): 1 / (1+b^2 t^2)$ \quad Коши $K(a): e^{-a|t|}$
\end{itemize}

\section*{Моменты и Свойства распределений}
Фундаментальное тождество: $\E[X^2] = \Var X + (\E X)^2$. 
\textbf{Симметрия модулей:} Для $X \sim \N(0, \sigma^2)$ или Лапласа $\E[X^{\text{нечет}}]=0$. Ковариация: $\Cov(|X|, X^2) = \E[|X|^3] - \E|X|\E[X^2]$. \textbf{Не нужно раскрывать модуль!} Вычисляется в 1 строку.
\textbf{Абсолютные моменты $\N(0, \sigma^2)$:}
$\E|X|^p = \sigma^p 2^{p/2} \Gamma((p+1)/2)/\sqrt{\pi}$.
$\E|X| = \sigma \sqrt{2/\pi}, \quad \E[X^2]=\sigma^2, \quad \mathbf{\E|X|^3 = 2\sigma^3 \sqrt{2/\pi}}, \quad \E[X^4]=3\sigma^4$.

\section*{Sanity Checks (Эвристики самопроверки)}
\begin{enumerate}
    \item \textbf{Размерность:} Если $[X]\!=\!\text{м}$, то $[\Var X]\!=\!\text{м}^2$. Для $g(X)\!=\!1/X^2$, дисперсия $\sigma_{as}^2 \sim [g']^2 \sigma^2 \!=\! \text{м}^{-6} \cdot \text{м}^2 = \text{м}^{-4}$. Нарушение = ошибка дифференцирования.
    \item \textbf{Квадратичная форма / Дисперсия:} $\sigma_{as}^2 \ge 0$, $\Var(Y|X) \le \Var(Y)$. Если $\sigma_{as}^2 = \nabla g^T \Sigma \nabla g < 0$ --- ошибка в знаках ковариации или Якобиане.
    \item \textbf{Закон сохранения массы:} Для смешанных строго $\int p_{ac} dx + \sum p_k \equiv 1$.
    \item \textbf{Тест константы (Вырождение):} $\E[Y^X]$ при $Y \!\sim\! U(1\!-\!\varepsilon, 1\!+\!\varepsilon)$. При $\varepsilon \!\to\! 0 \implies Y \!\to\! 1$. Формула обязана дать $\E[1^X] \!=\! 1$.
    \item \textbf{Оценка монотонности/Границы:} Если $Y \!\ge\! 1, X \!\ge\! 0 \implies Y^X \!\ge\! 1 \implies \E[Y^X] \!\ge\! 1$. В условном ожидании $\E[g(X,Y)|Z=z] = \varphi(z)$. Если $m \!\le\! g \!\le\! M$, то и $m \!\le\! \varphi(z) \!\le\! M$.
    \item \textbf{Теорема о двойном $\E$ (ФПМО):} Безусловное восстанавливается из условного: $\E[\E[g(X,Y)|Z]] \equiv \E[g(X,Y)]$.
    \item \textbf{Асимптотика экстремумов:} Если $Z \!=\! Y/X \to 0$ (при $X,Y \!>\! 0$), значит $Y \!\ll\! X$. Условная плотность $p_{X|Z}$ обязана концентрироваться у максимумов $X$.
\end{enumerate}

\section*{Аналитические методы вычисления интегралов и моментов}

\subsection*{1. Эйлеровы интегралы: Гамма- и Бета-функции}
Применяются для мгновенного вычисления моментов абсолютно непрерывных величин (экспоненциального, гамма-, бета- и хи-квадрат распределений), минуя многократное интегрирование по частям.

\textbf{Гамма-функция (обобщение факториала):}
\begin{gather*}
    \int_0^\infty x^{\alpha-1} e^{-\lambda x} \, dx = \frac{\Gamma(\alpha)}{\lambda^\alpha}, \quad \alpha > 0, \lambda > 0.
\end{gather*}
Для вычисления моментов $k$-го порядка гамма-распределения $X \sim \Gamma(\alpha, \lambda)$ плотность интегрируется с дополнительным множителем $x^k$, что эквивалентно сдвигу параметра формы: $\E[X^k] = \Gamma(\alpha+k) / (\lambda^k \Gamma(\alpha))$.

\textbf{Бета-функция (интегралы на конечном носителе):}
$
    \int_0^1 x^{\alpha-1} (1-x)^{\beta-1} \, dx = \operatorname{B}(\alpha, \beta) = \frac{\Gamma(\alpha)\Gamma(\beta)}{\Gamma(\alpha+\beta)}.
    $
Аналитический трюк для интегралов вида $\int_0^a x^p (a-x)^q \, dx$: замена $t = x/a$ сводит выражение к стандартной бета-функции с масштабным коэффициентом $a^{p+q+1}$.

\subsection*{2. Метод выделения полного квадрата (Гауссовский интеграл)}
Применяется при поиске математических ожиданий вида $\E[\exp(cX)]$ для $X \sim \N(\mu, \sigma^2)$ или вычислении нормировочных констант многомерных распределений.

Любой интеграл от экспоненты с квадратичным полиномом в показателе сводится к интегралу от плотности нормального распределения:
\begin{align*}
    \int_{-\infty}^\infty \exp\left( -ax^2 + bx + c \right) dx &= \exp\left(c + \frac{b^2}{4a}\right) \sqrt{\frac{\pi}{a}}, \quad a > 0.
\end{align*}
Для многомерного случая ($\mathbf{x} \in \mathbb{R}^n$, матрица $A$ положительно определена):
$
    \int_{\mathbb{R}^n} \exp\left( -\frac{1}{2}\mathbf{x}^T A \mathbf{x} + \mathbf{b}^T \mathbf{x} \right) d\mathbf{x} = \sqrt{\frac{(2\pi)^n}{\det A}} \exp\left( \frac{1}{2}\mathbf{b}^T A^{-1} \mathbf{b} \right).
    $

\subsection*{3. Характеристические функции и комплексная экспонента}
Используется для замены операции свертки плотностей алгебраическим умножением и дифференцирования вместо интегрирования при поиске начальных моментов.

\textbf{Генерация моментов:}
$
    \E[X^k] = \left. \frac{1}{i^k} \frac{d^k}{dt^k} \varphi_X(t) \right|_{t=0}.
        $
\textbf{Свертка независимых величин:} Если $Z = X + Y$ и $X \perp\!\!\perp Y$, вычисление $\int p_X(z-y)p_Y(y)dy$ заменяется тождеством $\varphi_Z(t) = \varphi_X(t)\varphi_Y(t)$ с последующим обратным преобразованием Фурье.

\textbf{Тригонометрические моменты (Формула Эйлера):}
При вычислении $\E[P(X)\cos(tX)]$ или $\E[P(X)\sin(tX)]$, где $P(X)$ — полином, осуществляется переход в комплексную плоскость:
\begin{align*}
    \E[\cos(tX)] &= \Re\left(\E[e^{itX}]\right) = \Re(\varphi_X(t)), \\
    \E[\sin(tX)] &= \Im\left(\E[e^{itX}]\right) = \Im(\varphi_X(t)).
\end{align*}
Для нормального распределения $\N(\mu, \sigma^2)$ результат получается немедленно, так как $\varphi_X(t) = \exp(i\mu t - \sigma^2 t^2 / 2)$.

\subsection*{4. Симметрия и нуллификация интегралов}
Применяется для мгновенного сокращения вычислений центральных моментов или ковариаций симметричных распределений.

Если случайная величина $X$ имеет плотность $p_X(x)$, симметричную относительно нуля ($p_X(x) = p_X(-x)$), и $\E|g(X)| < \infty$:
\begin{gather*}
    g(-x) = -g(x) \implies \E[g(X)] = \int_{-\infty}^\infty g(x)p_X(x)dx \equiv 0.
\end{gather*}
Следствия: все нечетные центральные моменты нормального распределения и распределения Лапласа равны нулю ($\E[X^{2k+1}] = 0$). Ортогональность полиномов разных четностей на симметричном интервале гарантирует, что $\Cov(X, X^2) = 0$, хотя величины функционально зависимы.
\subsection*{1. Условные распределения и ожидания}
\textbf{Условная плотность (Теорема Байеса для плотностей):}
Для абсолютно непрерывного вектора $(X, Y)$ условная плотность $X$ при условии $Y=y$ (где $p_Y(y) > 0$) определяется как:
\begin{gather*}
    p_{X|Y}(x|y) = \frac{p_{X,Y}(x, y)}{p_Y(y)} = \frac{p_{Y|X}(y|x) p_X(x)}{\int p_{Y|X}(y|x) p_X(x) dx}.
\end{gather*}

\textbf{Условное математическое ожидание (УМО):}
Является функцией от условия $\E[X|Y] = \psi(Y)$, минимизирующей среднеквадратичную ошибку прогноза $\E[(X - f(Y))^2]$. Вычисление в точке $Y=y$:
$
    \E[g(X) \mid Y=y] = \int_{\mathbb{R}} g(x) p_{X|Y}(x|y) \, dx.
    $

\textbf{Свойства УМО (Алгебра измеримости):}
    {Телескопическое (Башня):} $\quad \E[\E[X \mid Y]] = \E[X]$
    {Вынесение известного:} $\quad \E[h(Y) X \mid Y] = h(Y) \E[X \mid Y]$
    {Независимость:} $\quad X \perp\!\!\perp Y \implies \E[X \mid Y] = \E[X]$

\textbf{Закон полной дисперсии (Разложение дисперсии):}
\begin{gather*}
    \Var(X) = \E[\Var(X \mid Y)] + \Var(\E[X \mid Y]).
\end{gather*}

\subsection*{2. Преобразования случайных величин (Якобиан)}
\textbf{Одномерный случай (Диффеоморфизм):}
Пусть $Y = g(X)$, где $g$ — строго монотонная и дифференцируемая функция. Плотность распределения образа:
\begin{gather*}
    p_Y(y) = p_X(g^{-1}(y)) \cdot \left| \frac{d}{dy} g^{-1}(y) \right| \cdot I_{\{y \in \text{Im}(g)\}}.
\end{gather*}

\textbf{Многомерный случай (Векторное отображение):}
Для биективного преобразования $\mathbf{Y} = \mathbf{g}(\mathbf{X})$ совместная плотность определяется через определитель матрицы Якоби обратного отображения $\mathbf{x} = \mathbf{h}(\mathbf{y}) = \mathbf{g}^{-1}(\mathbf{y})$:
\begin{gather*}
    p_{\mathbf{Y}}(\mathbf{y}) = p_{\mathbf{X}}(\mathbf{h}(\mathbf{y})) \cdot |J|, \quad J = \det \left( \frac{\partial x_i}{\partial y_j} \right).
\end{gather*}

\textbf{Формула свертки (Сумма независимых величин):}
Если $Z = X + Y$, где $X \perp\!\!\perp Y$, плотность суммы:
$
    p_Z(z) = (p_X * p_Y)(z) = \int_{\mathbb{R}} p_X(x) p_Y(z - x) \, dx.
$

\subsection*{3. Базовые моменты и Неравенства}
\textbf{Ковариация и корреляция:}
$
    \Cov(X, Y) = \E[XY] - \E[X]\E[Y]. \\
    \rho(X, Y) = \frac{\Cov(X, Y)}{\sqrt{\Var(X)\Var(Y)}} \in [-1, 1].
$
\textbf{Дисперсия линейной комбинации:}
$
    \Var\left(\sum_{i=1}^n a_i X_i\right) = \sum_{i=1}^n a_i^2 \Var(X_i) + 2 \sum_{1 \le i < j \le n} a_i a_j \Cov(X_i, X_j).
$
\textbf{Фундаментальные неравенства:}
    \text{Маркова:} $\quad \P(|X| \ge a) \le \frac{\E[|X|^p]}{a^p}, \quad a>0, p>0$

    \text{Чебышёва:} $\quad \P(|X - \E X| \ge \varepsilon) \le \frac{\Var(X)}{\varepsilon^2}$

    \text{Йенсена:} $\quad f(\E[X]) \le \E[f(X)], \quad (\text{для выпуклой } f)$

    \text{Коши-Буняковского:} $\quad |\E[XY]| \le \sqrt{\E[X^2]\E[Y^2]}$

\section*{Дискретные распределения}
$F(x) = \sum_{k \le x} \P(X=k)$. Тождество: $\E[X^2] = \Var(X) + (\E[X])^2$.

\begin{itemize}[leftmargin=*, nosep, itemsep=4pt]
    \item \textbf{Бернулли $\text{Be}(p)$:} \\
    $\P(X=k) = p^k (1-p)^{1-k}, \quad k \in \{0, 1\}$ \\
    $\E[X] = p$, \quad $\Var(X) = p(1-p)$, \quad $\E[X^2] = p$

    \item \textbf{Биномиальное $\text{Bin}(n, p)$:} \\
    $\P(X=k) = \binom{n}{k} p^k (1-p)^{n-k}, \quad k = 0 \dots n$ \\
    $\E[X] = np$, \quad $\Var(X) = np(1-p)$

    \item \textbf{Пуассона $\text{Po}(\lambda)$:} \\
    $\P(X=k) = \frac{\lambda^k}{k!} e^{-\lambda}, \quad k = 0, 1, \dots$ \\
    $\E[X] = \lambda$, \quad $\Var(X) = \lambda$, \quad $\E[X^2] = \lambda(\lambda+1)$

    \item \textbf{Геометрическое $\text{Geom}(p)$ [от 1]:} \\
    $\P(X=k) = p(1-p)^{k-1}, \quad k = 1, 2, \dots$ \\
    $F(x) = 1 - (1-p)^{\lfloor x \rfloor}, \quad x \ge 1$ \\
    $\E[X] = \frac{1}{p}$, \quad $\Var(X) = \frac{1-p}{p^2}$, \quad $\E[X^2] = \frac{2-p}{p^2}$

    \item \textbf{Паскаля (Отриц. биномиальное) [до $r$]:} \\
    $\P(X=k) = \binom{k-1}{r-1} p^r (1-p)^{k-r}, \quad k \ge r$ \\
    $\E[X] = \frac{r}{p}$, \quad $\Var(X) = \frac{r(1-p)}{p^2}$
\end{itemize}

\section*{Непрерывные распределения}

\begin{itemize}[leftmargin=*, nosep, itemsep=5pt]
    \item \textbf{Равномерное $U(a,b)$:} \\
    $p(x) = \frac{1}{b-a}, \quad F(x) = \frac{x-a}{b-a}, \quad x \in [a, b]$ \\
    $\E[X] = \frac{a+b}{2}$, \quad $\Var(X) = \frac{(b-a)^2}{12}$

    \item \textbf{Экспоненциальное $\text{Exp}(\lambda)$:} \\
    $p(x) = \lambda e^{-\lambda x}, \quad F(x) = 1 - e^{-\lambda x}, \quad x \ge 0$ \\
    $\E[X] = \frac{1}{\lambda}$, \quad $\Var(X) = \frac{1}{\lambda^2}$, \quad $\E[X^2] = \frac{2}{\lambda^2}$

    \item \textbf{Нормальное $\mathcal{N}(\mu, \sigma^2)$:} \\
    $p(x) = \frac{1}{\sqrt{2\pi\sigma^2}} \exp\left(-\frac{(x-\mu)^2}{2\sigma^2}\right)$ \\
    $F(x) = \PhiFunc\left(\frac{x-\mu}{\sigma}\right)$ \\
    $\E[X] = \mu$, \quad $\Var(X) = \sigma^2$, \quad $\E[X^2] = \mu^2 + \sigma^2$

    \item \textbf{Лаплас $L(\mu, b)$:} \\
    $p(x) = \frac{1}{2b} \exp\left(-\frac{|x-\mu|}{b}\right)$ \\
    $\E[X] = \mu$, \quad $\Var(X) = 2b^2$, \quad $\E[X^2] = \mu^2 + 2b^2$

    \item \textbf{Коши $K(x_0, \gamma)$:} \\
    $p(x) = \frac{1}{\pi\gamma \left[1 + \left(\frac{x-x_0}{\gamma}\right)^2\right]}$ \\
    $F(x) = \frac{1}{\pi} \arctan\left(\frac{x-x_0}{\gamma}\right) + \frac{1}{2}$ \\
    \textit{Моменты $\E[X]$ и $\Var(X)$ не существуют.}

    \item \textbf{Гамма $\Gamma(\alpha, \lambda)$:} \\
    $p(x) = \frac{\lambda^\alpha}{\Gamma(\alpha)} x^{\alpha-1} e^{-\lambda x}, \quad x \ge 0$ \\
    $\E[X] = \frac{\alpha}{\lambda}$, \quad $\Var(X) = \frac{\alpha}{\lambda^2}$

    \item \textbf{Бета $\text{B}(\alpha, \beta)$:} \\
    $p(x) = \frac{1}{\operatorname{B}(\alpha, \beta)} x^{\alpha-1} (1-x)^{\beta-1}, \quad x \in [0,1]$ \\
    $\E[X] = \frac{\alpha}{\alpha+\beta}$, \quad $\Var(X) = \frac{\alpha\beta}{(\alpha+\beta)^2(\alpha+\beta+1)}$

    \item \textbf{Вейбулла $W(k, \lambda)$:} \\
    $p(x) = \frac{k}{\lambda} \left(\frac{x}{\lambda}\right)^{k-1} e^{-(x/\lambda)^k}, \quad x \ge 0$ \\
    $F(x) = 1 - e^{-(x/\lambda)^k}, \quad x \ge 0$ \\
    $\E[X] = \lambda \Gamma(1+1/k)$

    \item \textbf{Парето $\text{Pareto}(x_m, \alpha)$:} \\
    $p(x) = \frac{\alpha x_m^\alpha}{x^{\alpha+1}}, \quad F(x) = 1 - \left(\frac{x_m}{x}\right)^\alpha, \quad x \ge x_m$ \\
    $\E[X] = \frac{\alpha x_m}{\alpha-1}$, \quad $\Var(X) = \frac{\alpha x_m^2}{(\alpha-1)^2(\alpha-2)}$

    \item \textbf{Хи-квадрат $\chi^2_n$:} \\
    $p(x) = \frac{1}{2^{n/2}\Gamma(n/2)} x^{n/2-1} e^{-x/2}, \quad x \ge 0$ \\
    $\E[X] = n$, \quad $\Var(X) = 2n$, \quad $\E[X^2] = n(n+2)$

    \item \textbf{Стьюдент $t(n)$:} \\
    $p(x) = \frac{\Gamma(\frac{n+1}{2})}{\sqrt{n\pi}\Gamma(\frac{n}{2})} \left(1 + \frac{x^2}{n}\right)^{-\frac{n+1}{2}}$ \\
    $\E[X] = 0 \, (n>1)$, \quad $\Var(X) = \frac{n}{n-2} \, (n>2)$
\end{itemize}
\end{multicols*}
\end{document}
